\section{Kaggle lab: tune without touching the final test set}
\label{sec:kaggle-titanic-selection}

\begin{kaggleproject}{Titanic --- cross-validated logistic-regression selection}
\kagglesource{https://www.kaggle.com/datasets/yasserh/titanic-dataset}{Titanic Dataset}
\textbf{File.} \texttt{all\_datasets/titanic\_dataset/Titanic-Dataset.csv}\\
\textbf{Question.} Which regularization strength and class-weight policy perform best when tuning happens only inside the training data?\\
\textbf{Before running.} Draw the information boundary: outer train/test split first, then cross-validation only inside the training partition.

\begin{labmode}
This is the final guided Stage 3A lab. Trace the information boundary yourself before running the reference implementation; the important result is the selection procedure, not the winning hyperparameter.
\end{labmode}

\lstinputlisting[style=mlpython]{all_experiments/lab07-titanic_model_selection.py}

\textbf{Interpretation.} The training-CV score estimates the selection procedure inside development data. The final test score is a one-time estimate on untouched rows after the hyperparameters have already been chosen.

\begin{debugtask}
Move preprocessing outside the pipeline and fit it once on all rows before cross-validation. Then explain why that version leaks information even though the target was never explicitly passed to the preprocessor.
\end{debugtask}
\end{kaggleproject}

\begin{labdeliverable}
Submit the train/test information-boundary diagram, the cross-validation search summary, chosen hyperparameters, one final untouched-test result, and a short leakage explanation showing why fitting preprocessing before cross-validation would invalidate the selection estimate.
\end{labdeliverable}
